% ===== 原书 p.8 =====
\section{wrapped Floer 上同调}

由于 $Y$ 不是紧的，要讨论关于 $Y$ 的同调镜像对称，就不能用通常的 Fukaya 范畴，而必须用 wrapped Fukaya 范畴。它是辛上同调在 Lagrange 子流形情形下的类似物，其定义如下 [AS10]：设 $X^{\mathrm{in}}$ 是紧的带边流形，$\omega:=d\theta$ 是 $X^{\mathrm{in}}$ 上的辛形式，由 $\iota_{\vec Z}\omega=\theta$ 定义的 Liouville 向量场 $\vec Z$ 在边界 $\partial X^{\mathrm{in}}$ 上处处外向，此时称 $(X^{\mathrm{in}},\theta)$ 为 Liouville 域。这时 $(\partial X^{\mathrm{in}},\theta|_{\partial X^{\mathrm{in}}})$ 成为切触流形；沿 $-\vec Z$ 积分，该切触流形辛化的较小一半 $((0,1]\times\partial X^{\mathrm{in}},d(r\theta))$ 便成为 $X^{\mathrm{in}}$ 的领（collar）。注意辛化 $((0,\infty)\times\partial X^{\mathrm{in}},d(r\theta))$ 的 Liouville 向量场由 $\vec Z=r\partial_r$ 给出。把这个领与 $(\partial X^{\mathrm{in}},\theta|_{\partial X^{\mathrm{in}}})$ 辛化的较大一半 $([1,\infty)\times\partial X^{\mathrm{in}},d(r\theta))$ 粘合，所得的完全辛流形 $(X,d\theta)$ 称为 $(X^{\mathrm{in}},\theta)$ 的完备化；作为 Liouville 域的完备化所得的完全辛流形称为 Liouville 流形。Liouville 流形的 Lagrange 子流形 $L$，若 $\theta|_L$ 在紧集之外消失，则说 $L$ 在无穷远处是 Legendre 的。

固定 $X$ 上的一个函数 $H:X\to\R$，使其在紧集之外满足 $H(r,y)=r^2$。设 $\varphi_1:M\to M$ 是 $H$ 生成的时间 $1$ 的 Hamilton 流，则 $L\cap\varphi_1(L)$ 的元素称为 $L$ 的 Hamilton 弦（Hamiltonian chord）。这等价于说：$\gamma:[0,1]\to M$ 是 $H$ 生成的 Hamilton 向量场的积分曲线，且 $\gamma(0)\in L$、$\gamma(1)\in L$。
% [Tutor Note: 这一段及下一段中，原书把 Liouville 流形印作 $M$（「$\varphi_1:M\to M$」「$\gamma:[0,1]\to M$」「3 点付き円盤から $M$ への写像」），而本节的 Liouville 流形一开头即记为 $X$，$M$ 在此之前一直指格；按上下文三处均应作 $X$，译文照录原印]

给定 Hamilton 弦 $\gamma_+$ 与 $\gamma_-$，把从带状区域 $S:=\{s+t\sqrt{-1}\in\C\mid 0\le t\le 1\}$ 到 $X$ 的映射 $u:S\to X$ 中，满足 Floer 方程 (19) 与边界条件 $u(\partial S)\subset L$、$\lim_{s\to\pm\infty}u(s,t)=\gamma_\pm(t)$ 的个数记作 $n_{\gamma_-,\gamma_+}$。
\begin{equation}
\left(du-X_H\otimes dt\right)^{0,1}=0
\tag{19}
\end{equation}
\noindent 在由 Hamilton 弦生成的自由模 $CW^*(L)$ 上，用
\begin{equation}
\mathbf{m}_1([\gamma])=\sum_{\gamma'}n_{\gamma,\gamma'}[\gamma']
\tag{20}
\end{equation}
\noindent 引入微分，所得的复形称为 wrapped Floer 复形，其上同调称为 wrapped Floer 上同调。$CW^*(L)$ 的分次由 Maslov 指数定义。数出从带 $3$ 个标记点的圆盘到 $M$ 的、满足 Floer 方程与适当边界条件的映射的个数，就使 wrapped Floer 上同调带上分次环的结构。

在 $\overline{N}_{\C^\times}$ 上取对数极坐标 $(r_i,\theta_i)_{i=1}^{n-1}$，令自然的辛结构为 $\omega=\sum_{i=1}^{n-1}dr_i\wedge d\theta_i$。于是它成为以 $\theta=\sum_{i=1}^{n}r_id\theta_i$ 为 Liouville 形式的 Liouville 流形。
% [Tutor Note: Liouville 形式的求和上限在原图印作 $n$，而角坐标只有 $(r_i,\theta_i)_{i=1}^{n-1}$，多得一项 $r_nd\theta_n$；疑为 $n-1$ 之误，照录不改]
$\overline{L}_0:=\overline{N}_{\R^{>0}}$ 是满足 Legendre 条件 $\theta|_{\overline{L}_0}=0$ 的 $\overline{N}_{\C^\times}$ 的 Lagrange 子流形。取二次 Hamilton 函数 $H:=\frac{1}{2}|\bm r|^2=\frac{1}{2}(r_1^2+\cdots+r_{n-1}^2)$，则时间 $1$ 的 Hamilton 流为 $\varphi_1:(r_i,\theta_i)_{i=1}^{n-1}\mapsto(r_i,\theta_i+r_i)_{i=1}^{n-1}$，所以 $\overline{L}_0$ 的 Hamilton 弦与格点 $\bm n\in\overline{N}$ 一一对应。所有 Hamilton 弦的 Maslov 指数都是 $0$，从次数上看，Floer 复形的微分必为 $0$。wrapped Floer 上同调的积结构也不难算出，作为环同构于 $\overline{N}$ 的群环 $\C[\overline{N}]$。
% [Tutor Note: 原文此处印作「巻かれた Fleor コホモロジー」，Fleor 为 Floer 之误；译文按正字处理]

另一方面，$Y$ 是 $\overline{N}\times\C$ 的辛爆破中的开子集，且 $\overline{L}_0$ 与爆破中心不相交，因此 $L_0:=\overline{L}_0\times\R^{>0}$ 自然地确定 $Y$ 的 Lagrange 子流形。与 $\overline{N}\times\C$ 相比，$Y$ 的情形更为复杂，直接计算 $L_0$ 的 wrapped Floer 上同调并不容易；但当 $Y$ 的维数 $n$ 不超过 $3$ 时，仔细选取 Hamilton 函数，便可定义 wrapped Floer 上同调的一个变体 $\mathrm{HW}^*_{\mathrm{ad}}(L_0)$，使它满足下述性质：

% ===== 原书 p.9 =====
\begin{theorem}[{[CPUa, Theorem 1.2]}]
存在 $\C$ 代数的同构 $\mathrm{HW}^*_{\mathrm{ad}}(L_0)\cong H^0\left(\mathcal O_{\check Y}\right)$。
\end{theorem}

定理 5.1 的证明如下：在等式的左右两边都适当地取线性空间的基 $\{p_{\bm n,i}\}_{(\bm n,i)\in\overline{N}\times\Z}$，证明积结构可表为
\begin{equation}
p_{\bm n,i}\cdot p_{\bm n',i'}=\sum_{j=0}^{\ell(\bm n,\bm n')}\binom{\ell(\bm n,\bm n')}{j}p_{\bm n+\bm n',i+i'+j}
\tag{21}
\end{equation}
\noindent 即得。这里 $\ell:\overline{N}\times\overline{N}\to\overline{N}$ 是某个由组合方式定义的函数，它给出连接 Hamilton 弦的 Floer 方程的解与 $Z$ 的交点数。在 $\mathrm{HW}^*_{\mathrm{ad}}(L_0)$ 中 (21) 成立，由 [Pas14, Proposition 4.4] 推出；在 $H^0\left(\mathcal O_{\check Y}\right)$ 中 (21) 成立，则是简单的组合论。
% [Tutor Note: $\ell$ 的值域在原图印作 $\overline{N}$，但 (21) 把它用作二项式系数的上标，实际用到的是整数值（交点数）；照录不改]

作为环面 Calabi--Yau 流形的例子，有二维 $A_n$ 奇点的极小奇点解消、三维普通二重点（即所谓 conifold）的 crepant 解消；在这些情形下，可以比定理 5.1 更进一步，证明适当定义的 Fukaya 范畴与镜像的相干层的导出范畴等价。$A_n$ 奇点的极小奇点解消 $X$ 对应于这样的扇 $\Sigma$：其一维锥的生成元由 $v_0=(0,1),v_1=(1,1),\dots,v_n=(n,1)\in N\cong\Z^2$ 给出，此时的扇多面体为 $\Delta=[0,n]$。对以 $\Delta$ 为 Newton 多面体的 Laurent 多项式 $f(z)=(z-\alpha_0)(z-\alpha_1)\cdots(z-\alpha_n)$，考虑仿射超曲面
\begin{equation}
Y=\left\{(z,u,v)\in\C^\times\times\C^2\mid uv=f(z)\right\}
\tag{22}
\end{equation}
\noindent 此时无需考虑辛爆破，直接把 $\C^\times\times\C^2$ 的标准 K\"ahler 形式限制到 $Y$ 上所得的辛形式，即可用来讨论。$Y$ 上由正则体积形式
\begin{equation}
\Omega=\operatorname{Res}\left[\frac{dz}{z}\wedge\frac{du\wedge dv}{uv-f(z)}\right]=\frac{dz}{z}\wedge\frac{du}{u}
\tag{23}
\end{equation}
\noindent 赋以殆 Calabi--Yau 结构。映射
\begin{equation}
\pi_Y:Y\to B:=\R^2,\qquad (z,u,v)\mapsto\left(\log|z|,\ \frac{1}{2}\left(|u|^2-|v|^2\right)\right)
\tag{24}
\end{equation}
\noindent 是特殊 Lagrange 环面纤维化，其判别式集合由
\begin{equation}
\{(\log|\alpha_0|,0),\dots,(\log|\alpha_n|,0)\}\subset B
\tag{25}
\end{equation}
\noindent 给出。这些临界点称为焦点-焦点奇点（focus-focus singularity），经超 K\"ahler 旋转后，对应于小平（Kodaira）记号下称为 $I_1$ 型的最基本的椭圆纤维化的奇点（环面的一个闭链（cycle）塌缩，成为带普通二重点的有理曲线）。

$Y$ 的镜像 $\check Y$，是从 $X$ 中除去反典范因子 $\check D:=\chi_{(0,1)}^{-1}(1)$ 所得的流形。这里，$\chi_{(0,1)}$ 是与 $(0,1)\in M:=\Hom(N,\Z)$ 对应的特征标所伴随的 $X$ 上的函数。由于典范丛平凡，函数的零点给出反典范因子（注意这一点）。设 $E_1,\dots,E_n$ 是奇点解消的例外因子，则存在阿贝尔群的同构 $\operatorname{Pic}\check Y\to\Z^n,\ \mathcal L\mapsto\left(\deg\left(\mathcal L|_{E_i}\right)\right)_{i=1}^{n}$。经由这个同构，把对应于 $d\in\Z^n$ 的直线丛记作 $\mathcal L_d$。

% ===== 原书 p.10 =====
为简单起见，假设对任意 $i$，$\alpha_i$ 都是负实数，并令 $\Gamma=\{\alpha_0,\dots,\alpha_n\}$。对 $\C^\times\setminus\Gamma$ 上的道路 $\gamma:\R\to\C^\times\setminus\Gamma$，若 $|\gamma|:\R\to\R,\ t\mapsto|\gamma(t)|$ 是单调递增的双射，则称 $w(\gamma)=(w_1(\gamma),\dots,w_n(\gamma))$ 为其卷数；
% [Tutor Note: 原页作「単調増加な全単射」（500dpi 裁图核实：「全」在行末、「単射」在次行行首）。但 $|\gamma|\ge 0$，$\R\to\R$ 上不可能为满射，疑原书笔误；译文照录「双射」不改]
其中 $w_i(\gamma)$ 是 $\gamma$ 与闭区间 $[\alpha_{i-1},\alpha_i]$ 的交点数。对这个 $\gamma$，用以二次曲线为纤维的纤维化 $Y\to\C^\times,\ (z,u,v)\mapsto z$ 的纤维的适当 Lagrange 子流形，借 $Y$ 的辛形式作平行移动，便可作出 $\pi_Y$ 的 Lagrange 截面，其 SYZ 变换由 $\mathcal L_{-w(\gamma)}$ 给出 [CU13, Theorem 1.1]。进而，把这样得到的 Lagrange 子流形所生成的 wrapped Fukaya 范畴记作 $\mathcal W$，则有范畴等价
\begin{equation}
D^b\mathcal W\cong D^b\mathrm{coh}\,\check Y
\tag{26}
\end{equation}
\noindent [CU13, Theorem 1.2]。这些结果可以推广为由
\begin{equation}
Y:=\left\{(z,u_1,v_1,u_2,v_2)\in\C^\times\times\C^4\mid u_1v_1=z-a,\ u_2v_2=z-b\right\}
\tag{27}
\end{equation}
\noindent 定义的仿射流形与三维普通二重点的 crepant 解消 $X$ 中反典范因子 $\check D$ 的补集 $\check Y=X\setminus\check D$ 之间的镜像对称 [CPUb]。

% CHUNK c04 DONE pages p8–p10
